By a theorem of Beardwood, Halton and Hammersley, the length of a shortest tour through $n$ independent uniform random points in the unit square is asymptotic to $\beta\sqrt n$ for a constant $\beta$, which numerical estimates place near $0.712$. We prove $\beta<\MainBound$. The best earlier upper bound, $\beta\le0.8810$, is in a recent preprint of Dong and Cao. They cut the square into strips, take the points of each strip in blocks of six, and choose the order of visits within a block by a rule that is constant on the cells of a finite partition of the space of blocks. In each cell they bound the expected length of each segment of the path by the square root of its second moment. We change this last step: the length of a segment is a convex function of the positions of the points, and we bound its expectation by the Edmundson--Madansky inequality, a standard bound for convex functions, on a mesh much finer than the partition. For blocks of six points, $\CellsFive$ cells then give a better bound than their $1.6\cdot10^{13}$ cells. This makes blocks of \MainPoints{} points feasible, and these give the bound $\MainBound$. The proof ends in a computation in integer arithmetic, which three programs in different languages carry out with identical results.
